\documentclass[11pt,a4paper]{article}
\usepackage{fontspec,amsmath,amssymb,amsthm,microtype}
\setmainfont{lmroman10-regular.otf}[BoldFont=lmroman10-bold.otf,ItalicFont=lmroman10-italic.otf,BoldItalicFont=lmroman10-bolditalic.otf]
\usepackage{xcolor}
\usepackage[most]{tcolorbox}
\usepackage{mathrsfs}
\providecommand{\lt}{<}
\providecommand{\gt}{>}
\setmonofont{DejaVuSansMono.ttf}[BoldFont=DejaVuSansMono-Bold.ttf,ItalicFont=DejaVuSansMono-Oblique.ttf,BoldItalicFont=DejaVuSansMono-BoldOblique.ttf,Scale=0.86]
\newfontfamily\leanfallback{STIXTwoMath-Regular.otf}[Scale=0.86]
\definecolor{leanrule}{gray}{0.62}
\definecolor{leanback}{gray}{0.962}
\definecolor{leantitle}{gray}{0.22}
\definecolor{leanmuted}{gray}{0.38}
\newtcolorbox{leanbox}[1]{enhanced,breakable,colback=leanback,colframe=leanrule,
  boxrule=0.4pt,arc=0.6mm,left=2.4mm,right=2.4mm,top=1.8mm,bottom=1.8mm,
  before skip=0.8em,after skip=1em,fonttitle=\small\bfseries,coltitle=leantitle,
  colbacktitle=leanback,titlerule=0.3pt,title={#1}}
\newlength{\leancharwidth}
\newenvironment{leancode}{\par\smallskip\begingroup\ttfamily\footnotesize\raggedright
  \settowidth{\leancharwidth}{M}\setlength{\parskip}{0pt}\setlength{\parindent}{0pt}
  \frenchspacing}{\par\endgroup\smallskip}
\newcommand{\leanline}[2]{\hangindent=\dimexpr#1\leancharwidth+4\leancharwidth\relax
  \hangafter=1\noindent\hspace*{#1\leancharwidth}#2\par}
\newcommand{\leansource}[1]{\par{\footnotesize\color{leanmuted}#1}\par}
\newcommand{\leanlabel}[1]{\par\smallskip\noindent{\small\bfseries #1}\par}
\usepackage[margin=25mm]{geometry}
\usepackage[unicode,hidelinks]{hyperref}
\hypersetup{pdfauthor={LangevinLean Project}}
\setlength{\emergencystretch}{3em}
\setlength{\parskip}{0.4em}
\setlength{\parindent}{0pt}
\allowdisplaybreaks[2]
\raggedbottom
\title{\texorpdfstring{Real-range preservation by the actual Langevin operator}{Real-range preservation by the actual Langevin operator}}
\author{LangevinLean Project}
\date{}
\begin{document}\maketitle
Draft mathematical exposition; these notes may contain mistakes.


The canonical Brownian diffusion evaluations are actual measure-preserving observations. Their naturality identifies the complex transition on real observables with the actual real transition at every time. The real-range condition needed by the corrector is consequently derived.


\section{\texorpdfstring{Exact intertwining of the actual real and complex transitions}{Exact intertwining of the actual real and complex transitions}}


\textbf{Claim.} For every pointwise real isometry \(i\), every actual time \(t\) and real observable \(f\),

\[P_t^{\mathbb C}(if)=i(P_t^{\mathbb R}f).\]

\textbf{Proof.} Gibbs stationarity makes evaluation at zero and at \(t\) isometric pullbacks. Apply the exact compression embedding theorem to these two actual observations. This proves range preservation and applies in particular to the concrete orthonormal-coordinate inclusion, with no norm or dimension loss.

\begin{leanbox}{Lean: \texttt{instance\_\allowbreak{}at\_\allowbreak{}22}, \texttt{op\_\allowbreak{}mapLp}, \texttt{op\_\allowbreak{}mem\_\allowbreak{}real\_\allowbreak{}range} and 1 more}
\leansource{\texttt{instance\_\allowbreak{}at\_\allowbreak{}22} \quad WassersteinMixing/\allowbreak{}ComplexRealActual.lean, lines 22–22}
\begin{leancode}
\leanline{0}{local\ instance\ :\ InnerProductSpace\ ℝ\ H\ :=\ InnerProductSpace.rclikeToReal\ ℂ\ H}
\end{leancode}
\leansource{\texttt{op\_\allowbreak{}mapLp} \quad WassersteinMixing/\allowbreak{}ComplexRealActual.lean, lines 26–34}
\begin{leancode}
\leanline{0}{theorem\ op\_mapLp\ (i\ :\ R\ →ₗᵢ[ℝ]\ H)\ (P\ :\ Parameters)\ (V\ :\ E\ P.dimension\ →\ ℝ)}
\leanline{4}{(hV\ :\ PotentialAssumptions\ P\ V)\ (t\ :\ ℝ)\ (f\ :\ Lp\ R\ 2\ (ULA.gibbsMeasure\ V))\ :}
\leanline{4}{op\ ({\leanfallback 𝕜}\ :=\ ℂ)\ P\ V\ hV\ t\ (RealComplexEmbedding.mapLp\ i\ (ULA.gibbsMeasure\ V)\ f)\ =}
\leanline{6}{RealComplexEmbedding.mapLp\ i\ (ULA.gibbsMeasure\ V)}
\leanline{8}{(op\ ({\leanfallback 𝕜}\ :=\ ℝ)\ P\ V\ hV\ t\ f)}
\end{leancode}
\leansource{\texttt{op\_\allowbreak{}mem\_\allowbreak{}real\_\allowbreak{}range} \quad WassersteinMixing/\allowbreak{}ComplexRealActual.lean, lines 37–41}
\begin{leancode}
\leanline{0}{theorem\ op\_mem\_real\_range\ (i\ :\ R\ →ₗᵢ[ℝ]\ H)\ (P\ :\ Parameters)\ (V\ :\ E\ P.dimension\ →\ ℝ)}
\leanline{4}{(hV\ :\ PotentialAssumptions\ P\ V)\ (t\ :\ ℝ)\ (f\ :\ Lp\ R\ 2\ (ULA.gibbsMeasure\ V))\ :}
\leanline{4}{op\ ({\leanfallback 𝕜}\ :=\ ℂ)\ P\ V\ hV\ t\ (RealComplexEmbedding.mapLp\ i\ (ULA.gibbsMeasure\ V)\ f)\ ∈}
\leanline{6}{Set.range\ (RealComplexEmbedding.mapLp\ i\ (ULA.gibbsMeasure\ V))}
\end{leancode}
\leansource{\texttt{op\_\allowbreak{}coordinateLp} \quad WassersteinMixing/\allowbreak{}ComplexRealActual.lean, lines 45–52}
\begin{leancode}
\leanline{0}{theorem\ op\_coordinateLp\ (P\ :\ Parameters)\ (V\ :\ E\ P.dimension\ →\ ℝ)}
\leanline{4}{(hV\ :\ PotentialAssumptions\ P\ V)\ (t\ :\ ℝ)}
\leanline{4}{(f\ :\ Lp\ (E\ P.dimension)\ 2\ (ULA.gibbsMeasure\ V))\ :}
\leanline{4}{op\ ({\leanfallback 𝕜}\ :=\ ℂ)\ P\ V\ hV\ t}
\leanline{8}{(RealComplexEmbedding.coordinateLp\ (E\ P.dimension)\ (ULA.gibbsMeasure\ V)\ f)\ =}
\leanline{6}{RealComplexEmbedding.coordinateLp\ (E\ P.dimension)\ (ULA.gibbsMeasure\ V)}
\leanline{8}{(op\ ({\leanfallback 𝕜}\ :=\ ℝ)\ P\ V\ hV\ t\ f)}
\end{leancode}
\leanlabel{How the Lean proof follows the argument}
\texttt{op\_\allowbreak{}map\allowbreak{}Lp} specializes \texttt{compression\_\allowbreak{}map\allowbreak{}Lp} to the actual canonical state evaluations and their proved Gibbs laws. \texttt{op\_\allowbreak{}mem\_\allowbreak{}real\_\allowbreak{}range} gives the explicit range witness. \texttt{op\_\allowbreak{}coordinate\allowbreak{}Lp} uses the concrete coordinate embedding of the physical state space.
\end{leanbox}


Earlier results used here: \cite[Section~3.26.3]{LangevinLean2026ULASC}.

\section{\texorpdfstring{The corrector extraction is the actual real transition}{The corrector extraction is the actual real transition}}


\textbf{Claim.} A semigroup assembled from these actual canonical operators preserves the real coordinate image. Its extracted real transition \(i^*P_ti\) equals the actual real canonical transition itself.



\textbf{Proof.} The preceding intertwining identity puts every output in the embedded real range, which is exactly the projection-fixed condition. Apply the embedding adjoint to the same identity and use \(i^*i=I\) to identify the extracted transition.

\begin{leanbox}{Lean: \texttt{preservesReal\_\allowbreak{}of\_\allowbreak{}op\_\allowbreak{}eq}, \texttt{realOp\_\allowbreak{}eq\_\allowbreak{}op}}
\leansource{\texttt{preservesReal\_\allowbreak{}of\_\allowbreak{}op\_\allowbreak{}eq} \quad WassersteinMixing/\allowbreak{}ComplexRealActual.lean, lines 56–65}
\begin{leancode}
\leanline{0}{theorem\ preservesReal\_of\_op\_eq\ (i\ :\ R\ →ₗᵢ[ℝ]\ H)}
\leanline{4}{(P\ :\ Parameters)\ (V\ :\ E\ P.dimension\ →\ ℝ)\ (hV\ :\ PotentialAssumptions\ P\ V)}
\leanline{4}{(S\ :\ ReversibleSemigroup\ (Lp\ H\ 2\ (ULA.gibbsMeasure\ V)))}
\leanline{4}{(hS\ :\ ∀\ t,\ 0\ ≤\ t\ →\ S.op\ t\ =\ op\ ({\leanfallback 𝕜}\ :=\ ℂ)\ P\ V\ hV\ t)\ :}
\leanline{4}{ComplexRealCorrector.PreservesReal\ S}
\leanline{6}{(RealComplexEmbedding.mapLp\ i\ (ULA.gibbsMeasure\ V))}
\end{leancode}
\leansource{\texttt{realOp\_\allowbreak{}eq\_\allowbreak{}op} \quad WassersteinMixing/\allowbreak{}ComplexRealActual.lean, lines 69–79}
\begin{leancode}
\leanline{0}{theorem\ realOp\_eq\_op\ (i\ :\ R\ →ₗᵢ[ℝ]\ H)}
\leanline{4}{(P\ :\ Parameters)\ (V\ :\ E\ P.dimension\ →\ ℝ)\ (hV\ :\ PotentialAssumptions\ P\ V)}
\leanline{4}{(S\ :\ ReversibleSemigroup\ (Lp\ H\ 2\ (ULA.gibbsMeasure\ V)))}
\leanline{4}{(t\ :\ ℝ)\ (hS\ :\ S.op\ t\ =\ op\ ({\leanfallback 𝕜}\ :=\ ℂ)\ P\ V\ hV\ t)\ :}
\leanline{4}{ComplexRealCorrector.realOp\ S\ (RealComplexEmbedding.mapLp\ i\ (ULA.gibbsMeasure\ V))\ t\ =}
\leanline{6}{op\ ({\leanfallback 𝕜}\ :=\ ℝ)\ P\ V\ hV\ t}
\end{leancode}
\leanlabel{How the Lean proof follows the argument}
\texttt{preserves\allowbreak{}Real\_\allowbreak{}of\_\allowbreak{}op\_\allowbreak{}eq} rewrites the supplied operator identification and applies the proved actual range theorem. \texttt{real\allowbreak{}Op\_\allowbreak{}eq\_\allowbreak{}op} rewrites through \texttt{op\_\allowbreak{}map\allowbreak{}Lp} and cancels the isometric embedding with \texttt{adjoint\_\allowbreak{}embed}.
\end{leanbox}


Earlier results used here: \cite[Section~3.26.3]{LangevinLean2026ULASC}.

\begin{thebibliography}{ULA-LSI}
\bibitem[ULA-SC]{LangevinLean2026ULASC}
LangevinLean Project. \emph{Unadjusted Langevin under Strong Convexity}. LL-01, 2026. Draft report, content revision d34ead06987c. \url{https://odakyildiz.github.io/LangevinLean/docs/pdf/UnadjustedLangevin/FullProof.pdf}.
\end{thebibliography}

\end{document}
